\chapter{Bob Schutz}

The career of Bob Schutz was focused on foundational research in computational modeling of satellite trajectories and the Earth's gravity field. He had studied under Byron Tapley, who had recognized early on that accurately determining spacecraft trajectories in a non-uniform gravitational field would be a primary challenge of the space age, and had established the orbital mechanics program in 1961 and oversaw the integration of space science into the core engineering curriculum during the sixties and seventies. Schutz became one of Tapley’s earliest coworkers and was deeply involved with implementation of their research on the UT Computation Center 6600, Cyber, and Cray. The 6600 arrived in 1966, while he was a grad student, and he certainly ranked as an expert user of the 60-bit systems at UT, with roughly twenty years of experience from the arrival of the 6600, through the Cyber era and the Cray in the mid eighties. He developed a special geopotential formulation and other algorithmic innovations to take advantage of the computing hardware for specific improvements in applied gravity field modeling. This was hands-on research involving real-world data and the intricacies of very special computing hardware, stretching far back into the era of punchcards and magnetic tapes.

In the late sixties Schutz and Tapley attended a series of NASA workshops in Williamstown, Massachusetts. The 1967 meeting is often cited as the origin of the concepts that eventually led to the modern gravity mapping missions, and Tapley himself has noted that the need for a gravity-studying mission was defined during the 1967 session. During these deliberations, researchers including Tapley, who would later serve as the Principal Investigator for the Gravity Recovery and Climate Experiment, recognized that understanding dynamic terrestrial processes required unprecedented precision in orbital measurements. The theoretical foundations established at Williamstown directly anticipated the evolution of satellite-to-satellite tracking and gravity gradiometry, ultimately paving the way for transformative modern missions such as Grace, Goce, and Grace Follow-on, which have revolutionized our ability to monitor global water distribution, ice sheet mass balance, and ocean circulation.

Schutz is most prominently linked to the 1969 workshop and report The Terrestrial Environment: Solid Earth and Ocean Physics, which he has described as a visionary moment for the discipline. These workshops defined the scientific need for dedicated satellites designed for use in modelling and mapping the Earth's gravity field, and set the path forward for their work. The combination of  specialized satellites, radio and laser tracking stations, and computational modelling and simulation, all taking concrete form during the seventies, were the foundation for all that was to follow.

The 1969 report set specific goals that drove subsequent research efforts. For example, it called for measuring spatial geoid variations and mean sea level to within 10 cm, tracking geoid time variations to within 5 cm, and measuring relative crustal motion to an accuracy of 10 cm or better. To achieve these targets, the workshop explicitly recommended the development of drag-free, low-altitude satellites flying below 250 km, which would be equipped with radar altimeters. Crucially, the report proposed tracking these low-altitude satellites via satellite-to-satellite range-rate links from higher, stable relay satellites. These specific recommendations laid the theoretical groundwork for both the high-low and low-low satellite-to-satellite tracking configurations that materialized in later decades. By 1972, NASA had formally incorporated these Williamstown concepts into its Earth and Ocean Physics Applications Program, marking the moment when dedicated satellite gravimetry truly took shape.

To provide the higher, stable relay satellites recommended at Williamstown, the planning introduced the Geopause concept. Designed to orbit at an altitude of 30,000 km, it would be positioned far above the Earth's complex near-surface gravitational noise. At this height, the satellite would be perturbed by only six or seven low-degree gravity coefficients, allowing it to act as an exceptionally stable, drag-free anchor in space. From this vantage point, it could track low-altitude satellites via range-rate links to effectively isolate high-frequency gravity anomalies. To execute the low-low tracking configuration, the plan detailed the Gravsat mission concept, calling for a pair of low-orbiting satellites flying at an altitude of approximately 200 km, continuously measuring the changes in their inter-satellite range-rate to map high-resolution spatial gravity anomalies.

While it took decades for these dedicated gravity mapping constellations to be fully funded and engineered, the space community immediately began testing satellite-to-satellite tracking theories using available assets. The very first practical application of high-low tracking actually occurred during the Apollo program, when Earth-based tracking of Apollo lunar orbiters was used to successfully map lunar mass concentrations. In Earth orbit, early high-low tracking experiments utilized a relay satellite in geosynchronous orbit to track both the Apollo-Soyuz mission and a Geos spacecraft.

The roadmap outlined an integrated constellation of specialized spacecraft, explicitly planning for missions like the geodetic reference sphere Lageos and the Seasat mission for ocean dynamics. To provide the highly stable, unperturbed coordinate references required by the program, specialized passive satellites covered entirely in laser retroreflectors were launched into orbit. The French launched Starlette in 1975, and NASA launched Lageos in 1976. These missions provided the precise tracking targets necessary to study geopotential variations, tectonic plate motion, and Earth orientation parameters.

To handle the complex data generated by these new satellites and precision global tracking networks, the UT Austin team developed the University of Texas Orbit Processor software in 1974. Utopia became the computational workhorse for modeling the Earth’s gravity field and processing the new high-precision satellite laser ranging data to validate orbits with centimeter accuracy.

A recommendation of the 1969 report was the development of drag-free, low-altitude satellites flying below 250 km. In 1972 this ambitious engineering goal was first validated in space by the Navy's Triad mission, a member of the Transit navigation satellite family. Equipped with the Discos disturbance compensation system, Triad successfully demonstrated drag-free performance, proving that non-gravitational forces like atmospheric drag could be physically decoupled from orbital trajectories, proving that the ambitious low-altitude gravity mapping recommendations of the report were mechanically achievable.

Discos, developed collaboratively by Stanford University and the Johns Hopkins Applied Physics Laboratory, was a masterpiece of precision engineering based on a 1964 theoretical concept by Benjamin Lange. The heart of the system was its proof mass, a 22-millimeter sphere made of a highly dense gold-platinum alloy floating freely inside a 40-millimeter cavity. Because it was completely enclosed, the proof mass was shielded from all atmospheric drag and solar radiation pressure. When external forces pushed against the outer satellite, capacitive sensors detected the shifting gap between the cavity wall and the floating sphere. A control system then fired cold-gas thrusters to accelerate the spacecraft, forcing it to continuously chase the proof mass without ever touching it. In effect, the satellite was forced to fly along the proof mass's pure geodesic path.

Shielding the proof mass from the environment created the new problem of the gravitational mass of the satellite itself attracting the proof mass, so that any mass asymmetry in the satellite could ruin the drag-free state. To handle this situation, the satellite was designed as a tri-body structure. The central Discos unit was separated from the heavy power and main electronics units by 2.7-meter extendible booms, moving those masses far enough away to minimize their gravitational pull. Even the consumption of thruster fuel threatened to shift the satellite's center of mass and alter its gravitational pull on the proof mass. Researchers solved this by storing the propellant in two donut-shaped tanks placed symmetrically around the proof mass. As the gas depleted, its mass gradient relative to the center of the cavity remained at zero. 

The system worked flawlessly in orbit, successfully reducing non-gravitational effects to negligible levels. For the Navy's operational Transit navigation system, this meant orbital ephemeris predictions could be extended from twelve or sixteen hours out to a week or more without losing accuracy. More importantly for space geodesy, Triad proved that satellites could operate at the extremely low altitudes necessary to measure high-degree, short-wavelength gravity anomalies. Atmospheric drag at these low altitudes had been considered an overwhelming challenge, and the pioneering validation of drag-free metrology directly enabled later generations of gravitational observatories, including the European Space Agency's ultra-low Goce mission decades later.

Throughout the sixties and seventies, an important focus of Schutz and Tapley’s research was the Navy’s Transit navigational system and the quality of the orbital information being uploaded to the Transit satellites. Their work aimed to improve the system's accuracy and extend it for use in scientific and geodetic applications. This involved developing models and statistical filtering methods to process noisy observations from ground stations, moving away from simple models and towards those that accounted for relativistic effects and complex atmospheric drag. This era of computational modeling, filtering, and precision tracking culminated in the establishment of the Center for Space Research at UT Austin in 1981, with Tapley as founding director.

Their work explored adaptive data filtering procedures aimed at simultaneously improving both satellite orbit models and ground station positional accuracies. To handle the computational load of operational tracking for these dynamic systems, they developed efficient processors for adaptive state noise and dynamic model compensation. These formulations allowed their orbit determination software to adapt to unmodeled maneuvers and complex atmospheric drag variations without diverging or requiring manual tuning. In the sixties, this work relied on using doppler tracking data to evaluate radial orbit errors due to atmospheric interference and ionospheric delays. These were major noise sources and drove the shift toward satellite laser ranging in the seventies, an entirely new technique and a specialty of CSR and UT Austin. The core objective was to separate environmental propagation delays from the actual physical motions of the satellites, as this was a critical step in improving orbit predictability and precision.

The evolution of gravity mapping using satellites and the overall trajectory of the research at CSR began with a fortuitous discovery in the early sixties, driven by the urgent needs of Polaris missile submarine operations. Within a few years, a full-scale system of satellites and ground stations was deployed, regardless of cost, providing navigation fixes far beyond the reach of land-based beacons. The accuracy of these fixes turned out to be fundamentally constrained by knowledge of the Earth’s irregularities, geologic motions, and subsurface complexities. The entire paradigm was born days after the launch of Sputnik in 1957, when two physicists at the Applied Physics Laboratory, William Guier and George Weiffenbach, realized that they could determine the Soviet satellite's orbit by analyzing the doppler shift of its radio signals. Their manager Frank McClure recognized that this physical relationship could be inverted. If a satellite's orbit was already known and broadcasted, a navigator on Earth could measure the doppler shift to compute their exact geographic coordinates. This was exactly what the Navy needed for its Polaris submarine force. To maintain the concealment required for a credible nuclear deterrent, these submarines relied on inertial navigation systems which suffered from cumulative mechanical drift over prolonged deployments. The Transit satellite system allowed a submarine to briefly raise a small antenna to periscope depth and acquire a highly precise, passive position fix in any weather for resetting its inertial system.

The Navy’s plans quickly collided with the complexities of real-world experience. The Earth's non-spherical shape and uneven mass distribution continuously perturbed the Transit satellites' low-Earth orbits, degrading the accuracy of the orbital ephemerides transmitted to the submarines. Furthermore, the guidance systems for the ballistic missiles themselves were sensitive to these gravitational irregularities. Simulations showed that unmodeled crustal anomalies could introduce targeting errors exceeding two-hundred meters. Early precomputed implicit targeting models could not maintain the accuracy required, forcing a transition to explicit in-flight guidance algorithms that incorporated real-time gravity models. In short, to use a submarine-based missile accurately, the Navy first had to model the Earth's internal mass distribution.
This urgent strategic need catalyzed the field of precision orbit determination. Astrodynamicists realized that to know exactly where a satellite was, they had to abandon simplified analytical theories and use high-performance computers to numerically integrate the complex differential equations of motion, accounting for every physical force acting on the spacecraft. This led to early computational milestones like APL's orbit improvement program in 1959, and subsequently, advanced software like Utopia at UT Austin. The mathematical and engineering frameworks developed to ensure the accuracy of submarine-launched ballistic missiles became the primary tools for measuring the Earth's detailed physical parameters. By applying these POD principles to advanced configurations, like the low-low satellite-to-satellite tracking envisioned at the 1969 Williamstown workshop, modern missions such as Grace now continuously measure the subtle gravitational pulls of shifting reservoirs and deep-ocean currents.

The problem in the early sixties was to provide periodic updates or corrections to the inertial navigation systems onboard Polaris submarines. Inertial or black-box navigation does not require external inputs, and so allows the submarines to remain concealed, but over time drifts and other errors build up. By acquiring external navigation fixes, the inertial errors can be kept below a threshold, with the tradeoff of reducing the submarine’s concealment during the fix. A number of methods were energetically developed for external fixes, but the leader was the Navy Navigation Satellite System, commonly known as Transit, the first operational satellite-based navigation system. The system relied on measuring the radio signals transmitted by Transit satellites. By analyzing how the signal frequency changes as it approaches and recedes it is possible to calculate the satellite's speed and closest approach. The submarine's onboard systems derived their exact two-dimensional position from the specific shape of the received doppler curve. As the satellite passed overhead, the exact moment the received frequency matched the broadcast frequency, the zero-doppler shift point, indicated the time of closest approach. Furthermore, the steepness or rapidity of the frequency change revealed the relative longitude between the submarine and the satellite. Because the Earth rotates beneath the satellite, the doppler curve also contained a slight asymmetry that allowed the system to determine whether the satellite was passing to the east or the west of the submarine.

Achieving the required accuracy meant overcoming several major physical and engineering hurdles. The ionosphere bends and delays radio waves, which would normally corrupt the doppler curve and ruin the position fix, so Transit satellites were designed to broadcast simultaneously on two distinct frequencies, 150 MHz and 400 MHz. Because the ionosphere's refractive effect is frequency-dependent, the submarine's receiver could mathematically compare the two signals and cancel out the atmospheric delay, ensuring a highly accurate fix. In the late fifties, processing this complex doppler curve required digital computers, which were at the time still typically room-sized machines. To bring this capability onboard submarines, engineers had to design a pioneering, highly influential military tactical computer. Introduced in 1961 by Ramo-Wooldridge, the TRW-130 was designated as the AN/UYK-1 by the Navy. It is historically famous as the first computer small enough to fit through a submarine hatch, partly thanks to specially rounded corners. It could process the satellite's orbital parameters and doppler data to generate a position fix in about fifteen minutes. While the system was designed for a constellation of evenly spaced satellites in polar orbits at an altitude of about 1,100 km, launch vehicle injection errors caused the orbital planes to precess over time. This led to a clustering phenomenon known as the streetcar effect, where several satellites would pass over in rapid succession, followed by long gaps of several hours where no satellites were above the horizon, forcing submarines to wait to safely reset their inertial systems. 

The system relied on the satellites broadcasting their own known orbit. Because these low-Earth orbits were constantly degraded by unmodeled atmospheric drag and Earth's uneven gravity, the Navy had to continually track the satellites from the ground and upload fresh orbital ephemerides and clock corrections twice per day. This limitation is what ultimately drove the continuous push for better gravitational models and advanced orbit determination techniques. The upload data was generated by the Transit system’s ground stations. This ground segment was the brain that kept the Transit system accurate. Tracking stations like the one operated by UT Austin at McMurdo Station Antarctica, along with a similar station in Thule Greenland, were strategic assets. Because the Transit satellites flew in polar orbits, these high-latitude ground stations were able to observe and track every pass of every Transit satellite. Every time a Transit satellite passed overhead, the stations recorded the doppler shift of its signal. During the seventies, Station 019 at McMurdo was typically staffed by electrical engineering graduate and undergraduate students from UT Austin, who operated the equipment out of a small quonset hut. The station measured the doppler shifts, transferred the raw tracking data onto 5-hole paper tape, and then transmitted it via commercial and military teleprinter networks back to the Satellite Control Center at the Applied Physics Laboratory in Maryland.

Once APL received this global tracking data, their central computers processed it using sophisticated software that calculated a trial doppler curve based on an initial estimated trajectory and then performed a recursive least-squares curve fit against the actual doppler data collected by the ground stations. This created a precise determination of the satellite's true orbit despite the complex environmental forces acting upon it. After calculating the updated orbital parameters, the Navy used one of four dedicated tracking and injection stations to upload the fresh orbit ephemeris and clock corrections to the satellites twice each day. Each satellite stored an updated navigation message in its onboard computer memory and broadcasted it back down to the concealed submarines as a phase-modulated signal on the 150-MHz and 400-MHz radio channels. Even though updates happened twice a day, the satellites usually carried enough orbital data to function for up to sixteen hours in case a ground station missed an upload window. The data uploaded from the ground consisted of a 1500-bit navigation message repeated every two minutes by the satellites. This message contained the general shape of the orbit, variable parameters describing the deviations caused by the Earth’s gravity anomalies and atmospheric drag, and timing corrections for small adjustments to the satellite's internal clock to stay synchronized with ground time. 

The accuracy of the system was built on, and completely reliant on, the quality of the Earth gravity models used. In the early days of Transit, the gravity models were so poor that drift in the predicted positions of the satellites was the biggest source of overall error in the navigational fixes onboard the submarines. As the system evolved over time however, the ground stations helped map the Earth's shape and gravity field more accurately, which in turn improved the satellite fixes. At some point there was a crossover, when the largest error source switched over and became the quality of the orbital information that was uploaded to the satellites, and this is where UT Austin directly entered the story. By recording the doppler shifts of every pass at ground stations, like UT’s Station 019 in Antarctica, researchers amassed a global dataset of orbital perturbations. The problem of computing and uploading highly accurate orbital information led Byron Tapley and Bob Schutz to recognize that the classical, analytical perturbation theories used at the time could not absorb the complex, non-linear reality of the near-Earth environment. To provide the accuracy the system demanded, they drove a shift toward computational modeling, simulation, and  numerical state propagation.
 
Instead of using simplified mathematical expansions, Tapley and Schutz leveraged the processing power of the UT Computation Center CDC 6600 to numerically integrate the full, non-linear differential equations of satellite motion. This allowed them to directly simulate complex, localized gravitational anomalies that analytical theories missed. While gravity was a major factor, atmospheric drag in low-Earth orbit was highly unpredictable due to fluctuations in solar activity and atmospheric density, and the drag signature was often indistinguishable from gravitational perturbations. The UT Austin group pioneered advanced statistical filtering techniques, allowing the software to continually estimate and compensate for dynamic model errors. In 1974 they consolidated these advanced physical force models, numerical integration techniques, and statistical estimation algorithms into Utopia, the University of Texas Orbit Processor. Utopia became the computational engine capable of digesting vast amounts of noisy radiometric doppler data and separating atmospheric drag from true gravitational variations. The mathematical and software frameworks that UT Austin established for cleaning up the Transit ephemeris uploads became the foundation for subsequent gravity modeling.

\section{Geodesy And Spherical Harmonics}

Geodesy is the study of the geometry, gravity, and spatial orientation of the Earth, and while treating the Earth as a sphere with uniform mass distribution is useful for preliminary analysis, this over-simplified model breaks down quickly because of real-world oblateness, equatorial bulges, localized mass concentrations, and higher-order anomalies. If these irregularities are ignored, the inaccuracies in the gravity model accumulate, leading to significant orbit prediction errors. During the development of the Transit system, early orbital models produced trajectory prediction errors of up to two kilometers over a 24-hour period. To achieve the system's operational requirement of providing submarines with positional accuracy of better than 0.1 nautical miles, researchers had to greatly improve the working models of the Earth's non-uniform gravitational field.

The Earth’s gravitational field conserves energy and has no local rotational forces. Mathematically, this specific physical reality is summarized and enforced by Laplace's equation. Any function that satisfies Laplace's equation is known as a harmonic function because the solutions are deeply connected to the equations that describe vibrating strings and musical harmonics. A harmonic function can never have a peak or a valley inside its domain, unless it's completely flat. Any highs or lows must occur on the boundaries. Imagine a stretched rubber sheet. It can curve, but it won't have any random bubbles or dips in the middle. When Laplace's equation is formulated and solved using spherical coordinates, the solution can be separated into radial and angular components, and the angular components are discrete “musical notes” over the surface of the sphere.

By formulating the gravitational field as a series expansion of spherical harmonics, researchers can systematically isolate the dominant point-mass potential from the additional contributions that encode a body's deviations from spherical symmetry. This spherical harmonic expansion describes the physical geometry of the Earth's gravity field by the degree and order of the expansion, and these divide naturally into distinct geometric patterns. The zonal harmonics depend exclusively on latitude and describe variations that are rotationally symmetric from the North Pole to the South Pole, such as the Earth's equatorial bulge and its pear-shaped north-south asymmetry. The sectorial harmonics vary longitudinally and divide the Earth into alternating positive and negative gravitational sectors, much like the slices of an orange. The tesseral harmonics vary in both latitude and longitude, forming a complex checkerboard pattern of mass distribution anomalies across the globe. The introduction of spherical harmonic models provided an expressive framework that was essential for understanding dynamic orbital processes.

There are some similarities with the finite elements used in subsurface modeling. They have a shared philosophical and structural lineage, based on the discretization of continuous phenomena into manageable mathematical components, and the pursuit of high-fidelity approximation through the summation of simpler, localized, or orthogonal functions. Both methods conclude that the most effective way to understand a complex, irregularly shaped reality is to decompose it into a hierarchy of simpler, mathematically rigorous parts, whether those parts are the physical bricks of a subsurface mesh or the harmonic ripples of a planetary gravity map. In the context of subsurface modeling, finite elements address the inherent heterogeneity of the Earth’s lithosphere by partitioning a continuous volume into a mesh of discrete geometric shapes, such as tetrahedra or hexahedra. This patchwork approach allows researchers to assign unique physical properties such as porosity, thermal conductivity, or elastic moduli to specific regions, thereby capturing the non-linear behavior of seismic waves or fluid flow within a fragmented geological medium. This methodology mirrors the operational logic of spherical harmonics used in gravity modeling, despite the latter’s reliance on global rather than local basis functions. Where finite elements utilize a localized mesh to resolve small-scale variations in rock density, spherical harmonics employ a series of undulating, overlapping waves to reconstruct the Earth’s total mass distribution. 

This relationship of finite elements and spherical harmonics can be conceptualized through the analogy of a digital image. Finite elements function similarly to a rasterized grid of pixels, where each cell holds specific data to form a local picture. Spherical harmonics operate like a compression scheme, where the image of the gravity field is built by layering broad, global patterns with increasingly fine, high-frequency overtones. Just as increasing the number of elements in a mesh improves the resolution of a fault line, increasing the degree and order of a harmonic expansion allows for the detection of smaller gravitational anomalies, such as those caused by mountain ranges or deep-seated mantle plumes. Spherical harmonics can also be conceptualized through the analogy of musical acoustics or the vibrations of a drum. Just as a complex musical note is composed of a fundamental frequency and a series of overtones, each adding specific texture and detail to the sound, the Earth’s gravitational field is reconstructed by layering various global tones. At the most basic level, the fundamental tone represents the gravity of a perfect sphere. However, because the Earth’s rotation causes it to bulge at the equator, a second layer, or overtone, is added to account for this elliptical distortion. As the model increases in complexity, finer and finer layers are added to account for massive mountain ranges, deep ocean trenches, and the dense metallic compositions of the tectonic plates.

In practice, there are many challenges involved in modeling the Earth’s gravity using spherical harmonics. The coefficients in the harmonic series are notoriously difficult to determine experimentally with sufficient accuracy. Researchers require a uniformly distributed, worldwide tracking network to gather dense data. Any large gaps in data coverage present a serious difficulty for accurately evaluating the coefficients. One challenge is non-zonal harmonics. Unlike zonal harmonics, which are symmetric around the Earth's axis and cause long-term, secular drifts in an orbit, non-zonal harmonics are longitude-dependent. Because the Earth rotates underneath the satellite, these non-zonal gravity anomalies exert a periodic push and pull on the spacecraft. The fundamental period of these forces is roughly the time it takes the Earth to make one complete revolution beneath the orbit. Because these forces reverse sign so rapidly, the satellite simply does not have enough time to accelerate into large positional deviations before the gravity force pushes it back the other way. As a result, the physical oscillations in the satellite's position are tiny, typically only a few tenths of a mile.

To measure these rapid, small-amplitude oscillations, researchers need data that is extremely dense in time. Ideally, at least one set of tracking data for every single revolution the satellite makes around the Earth. If a tracking network has large geographical gaps, it causes gaps in the data. If those data gaps happen to synchronize with specific times of day, researchers essentially develop a blind spot and cannot accurately follow the full, continuous cycle of the orbital oscillations. Early optical tracking cameras could not meet these requirements because they were heavily dependent on clear weather and nighttime skies. Tranet was the first real solution, with a globally distributed network of ground stations that continuously tracked the radio doppler frequency shifts of the Transit satellites, immune to weather and daylight. At its peak, Tranet was gathering upwards of 100 sets of doppler data per satellite, per day.
Even with dense data, measuring the gravitational field requires observing how the field affects different orbital geometries. To mathematically separate and accurately evaluate the hundreds of different harmonic coefficients, researchers had to collect this dense data from multiple satellites operating at distinctly different orbital inclinations. For example, just to definitively separate the harmonics up to degree and order 8, researchers needed data from a minimum of four different orbital inclinations.

As models grew more complex, researchers attempted to map high-order terms, degree and order greater than 100, which represent very short-wavelength, localized gravity variations. However, tracking these micro-variations hit a physical wall due to the Earth's atmosphere. Radio signals traveling from space to a ground station must pass through the ionosphere, which stretches the signal wavelength, and the troposphere, which slows the signal down. While researchers developed brilliant workarounds, such as transmitting at both 150 and 400 MHz, to cancel out ionospheric dispersion, and using models to predict tropospheric weather delays, the atmosphere still leaves trace amounts of random noise in the tracking data. For very high-order gravity anomalies, the physical perturbation on the satellite is so minuscule that the atmospheric noise in the ground-based tracking signal is actually larger than the gravitational effect being measured. To overcome this atmospheric limitation, modern and future gravity mapping missions rely on satellite-to-satellite tracking or orbiting gravity gradiometers, allowing them to measure tiny gravitational perturbations while staying entirely above the Earth's noisy atmosphere.

In the early sixties, satellite geodesists initially assumed that the spherical harmonic expansion of the Earth's gravity potential would be a rapidly converging series. They believed they would only need to calculate a handful of low-degree terms to accurately model a near-Earth satellite's orbit. However, as tracking networks collected more data and researchers attempted to estimate more coefficients, they realized the convergence was much slower than anyone had anticipated. The first definitive proof of this slow convergence came from the discovery of orbital resonance. Transit 5BN-2 was found to have a massive, unexplained along-track oscillation of 130 meters that repeated every 60 hours. Transit 5E-1 showed a similar tracking oscillation with a 65-hour period. Researchers soon realized that these 60-hour and 65-hour beat periods were the result of the satellites entering a near-resonant orbit with the Earth's 13th-order spherical harmonics. Upon closer inspection of the data, they also found a smaller 40-hour beat period, which meant the satellites were simultaneously experiencing a secondary resonance with the Earth's 14th-order harmonics. The unexpectedly large amplitudes of these specific high-order coefficients proved that fine-scale gravitational variations exerted a significant force and could not be ignored. 

Further complicating the computation, harmonic coefficients naturally span multiple orders of magnitude, even among low-degree and low-order terms. This vast numerical spread makes it difficult to intuitively assess their relative physical contributions to the gravity field. To handle this, scientists often normalize the coefficients using a reference length scale, such as the Earth's equatorial radius, to make them dimensionless and easier to compare. Because the harmonic series converges so slowly, truncating a computational model too early leads to severe aliasing effects. When the orbit determination software tries to reconcile the satellite's actual tracking data with an incomplete gravity model, it misinterprets those unmodeled forces. The software essentially aliases the missing high-order forces into the solution, which artificially warps and corrupts the estimates of the lower-order coefficients that are included. Early results found that the unmodeled forces from the 13th, 14th, and 15th-degree resonant harmonics were actually the major contributor to errors in their estimates of the lower-order non-zonal coefficients, degree 10 and below. To achieve high accuracy, scientists had to map the gravity field to a much higher degree than originally planned just to prevent the ignored high-order terms from distorting the earlier, nominally simpler, parts of the model.

On top of all these physical, real-world challenges, the numerical computations involved in processing early spherical harmonics models of the Earth’s gravity field were daunting. This was exactly the field that Bob Schutz entered in the sixties, just as the CDC 6600 was arriving at UT Austin. Modeling and simulation on the 6600 was his expertise, and he developed several important innovations tied directly to the numerical code created for gravity modeling on UT Austin’s machine. It is fair to say that UT Austin was uniquely prepared for the gravity modeling challenges of the sixties and seventies purely because of its 60-bit floating-point oriented computer. 

\section{CDC 6600 And Cyber}

Computing spherical harmonics involves recursive equations and repeated differentiation of lower-order terms. As researchers attempt to model higher degrees of the gravitational field, this process becomes  computationally expensive and prone to numerical instability. The models require solving for hundreds or even thousands of highly correlated mathematical coefficients using heavy computational power. The reason the coefficients are so highly correlated is rooted in the physics of orbital mechanics. Gravitational force decays with distance, and at satellite altitudes the perturbation signals from high-degree gravity anomalies, features smaller than a continent, are highly attenuated. When this attenuation was combined with the sparse distribution of early ground tracking stations, it resulted in ill-conditioned normal matrices that severely strained the era’s mainframes. Historically, processing just one iteration of a gravity model could take up to fifty hours on early mainframe computers, and modern high-degree models require the use of advanced supercomputers.

To avoid having to mathematically invert the entire massive geopotential matrix at once, researchers looked for specific beat frequencies where a satellite's orbit resonated with the Earth's rotation. For example, a satellite orbiting at 12.82 cycles per day experiences a primary resonance with 13th-order gravity terms. Over a 5.5-day period, these resonant terms would dramatically amplify, causing an along-track perturbation of nearly 600 meters. Researchers used these resonant effects to isolate and solve for specific high-frequency coefficients, such as 27th-order terms, without triggering a 50-hour computational marathon. Early software systems like Celest used an iterative batch least-squares approach, meaning the computer attempted to process all tracking observations simultaneously to minimize squared residuals. To ease the computational burden and better handle dynamic errors like atmospheric drag, researchers developed sequential estimators like Prefer. Instead of loading the entire global tracking matrix, Prefer processed data in 120-second mini-batches using a Kalman filter. Crucially, as tracking stations or satellites passed out of view, the software dynamically dropped their associated parameters from the active mathematical state vector, which reduced memory usage.

The computing hardware used for gravity modeling and orbit determination stretched back to the early fifties. In September 1950, seeking a highly reliable and fast computer for complex ballistic equations, the Naval Bureau of Ordnance contracted IBM to build a custom machine. IBM developed the computer under a cost-plus fixed-fee contract, earning a nominal profit of exactly one dollar. Directed by Wallace Eckert and assembled at Columbia University's Watson Scientific Computing Laboratory, the machine cost approximately $2.5 million to build, roughly equivalent to $25 million today. The Norc was a room-sized, vacuum-tube-based machine that was widely considered the world's first supercomputer. It utilized an electrostatic memory system capable of storing 4,096 words, with auxiliary storage provided by up to eight magnetic tape units. IBM publicly demonstrated the Norc on December 2, 1954. Although initially designated for the Naval Ordnance Laboratory at White Oak, the Navy ultimately decided to send it to Dahlgren Virginia. In March 1955, the massive computer was disassembled, packed into seventeen trucks, and shipped to the Naval Ordnance Computation Center, where it became fully operational in July 1956. It operated 24 hours a day, seven days a week, supported by an expanded staff of 30 programmers and 30 maintenance personnel.

During its design, a major debate occurred regarding the computer's mathematical architecture. While some prominent scientists like Jon von Neumann preferred binary notation, IBM and the Navy committee opted for a decimal numbering system utilizing a 14-digit word length, because they believed it would make spotting human errors much easier. Dahlgren personnel also successfully lobbied IBM to include a self-checking feature using bit-count checks, which caught transfer errors with minimal hardware overhead. These design choices made the NORC stable, allowing it to spend more than 90\% of its time online. Operating at a speed of 15,000 operations per second, its developers noted that it calculated "like a human but completing it in a millionth of a second".

Running a single 6-degree-of-freedom trajectory simulation took 1.5 hours. This became a critical issue when the Navy needed to compute targeting data for the Polaris submarine-launched ballistic missile. Planners estimated that running the required trajectories would take 40 years of continuous computing on the Norc. To bypass this hardware limitation, Dahlgren researchers Russell Lyddane and Russell Cohen invented a workaround. Rather than simulating every possible path, they ran just enough 1.5-hour trajectories on the Norc to develop a precomputed grid interpolation process. This allowed the submarine's onboard fire control computers to quickly generate targeting data in real time using numerical interpolation functions, completely circumventing the mainframe's speed limits.

Patricia Smith wrote the Orki program using the Norc’s native machine code to compute the first definitive orbits for the early Transit navigation satellites. Paul Herget used the Norc to calculate the most precise measurements of Earth's orbit at the time. After completing a massive 9-hour computational run, he was so impressed that he named an asteroid The Norc. Dahlgren eventually supplemented the Norc with IBM 7090 mainframes. By the early sixties, it was superseded by the next generation of transistorized mainframes, such as the IBM 7030 Stretch, which could perform hundreds of thousands of operations per second.

A historic computing marathon occurred when researchers attempted to simultaneously evaluate 70 different geopotential coefficients, mapping the gravity field up to degree and order 8. To do this, the computer had to perform a massive least-squares fit that adjusted satellite orbit parameters and ground station locations against 1,662 separate sets of tracking data. Once the data was pre-processed and selected, this single overarching computation took approximately fifty hours to run on an IBM 7094 mainframe. To maintain numerical stability when inverting a dense, highly correlated matrix for 70 geopotential coefficients alongside the coordinates of 49 ground stations and various satellite orbital parameters, the calculations had to be performed in 72-bit double precision. The transistorized IBM 7094 possessed a physical core memory of only 32,768 words, which equated to roughly 144 kilobytes. Because the massive normal equation matrix far exceeded this capacity, programmers were forced to implement out-of-core block-partitioning algorithms. The software divided the dense matrix into smaller sub-matrices that were sequentially saved to and read from external IBM 729 magnetic tape drives. The mechanical seek, rewind, and read/write times of these tape units were relatively slow, meaning the central processing unit spent a huge fraction of those 50 hours simply waiting for data transfers to complete. 

Despite the grueling processing time, this specific computation was a success. Conducted by researchers at the Naval Weapons Laboratory led by Richard Anderle, it produced the model validating the use of the Tranet radio doppler tracking network for global geodesy, proving that ground stations could be located to an accuracy of 20 meters, which prompted the military to adopt doppler tracking as its primary geodetic reference. To handle the steadily growing matrix calculations, military laboratories had to continuously invent new generations of orbit determination software. They progressed from Orki on the Norc computer, to Astro on the transistorized IBM 7030 Stretch, and eventually to Celest. The transition to the IBM 360 mainframe in 1964 was a major milestone, allowing programmers to use standardized languages like Fortran and assembly to process the large global geodetic matrices.

As the tracking data grew denser, the models grew exponentially larger. By the early seventies, gravity models such as WGS-72 required estimating about 500 terms using IBM 7094 and 360 computers. Modern high-degree models, complete to degree and order 70, require solving for more than 7,000 highly correlated terms. To process this modern workload without taking weeks, researchers rely on parallel processing and advanced supercomputers, such as the Cray, which can complete a 7,000-term gravity solution in about two hours. 

Returning to the subject of computing spherical harmonics via recursive differentiation, researchers developed numerical innovations to overcome the severe numerical instability and computational expense of modeling higher-degree gravity anomalies. To eliminate the need for repeated differentiation, they replaced the standard formulation using Rodrigues' rotation formula. Instead, they switched to three-term algebraic equations so that by starting with a simple base value of 1, computers could calculate higher-order shapes using recursive formulas that held either the polynomial's degree or order fixed at each step. This approach required only basic algebraic operations and remained strictly numerically stable. 
Even with stable recursive math, traditional latitude-dependent spherical harmonic models suffered from mathematical singularities at the Earth’s poles. Schutz developed an approach based on a normalized version of the Pines formulation. While Pines had proposed using Cartesian coordinates to bypass polar singularities, there was still instability at high degrees due to large factorial terms. By introducing a scaling factor for the problematic terms, The new approach maintained complete numerical stability, allowing the geopotential field to be evaluated safely up to very high degrees. As global gravity models rapidly expanded in resolution, evaluating these complex functions remained a bottleneck. Schutz also helped derive highly optimized recursion relations that replaced complex transcendental functions with simple algebraic operations explicitly designed to exploit the parallel and vector pipeline execution capabilities of new Cray supercomputers.

To manage the demands of continuous data ingestion alongside high-degree spherical harmonic modeling, the UT Computation Center underwent several major hardware upgrades over the decades. Its CDC 6600 was uniquely suited for early orbit determination because it used 10 independent peripheral processing units to handle input and output operations, which freed the main processor to focus strictly on complex numerical calculations. Because the 6600 utilized a 60-bit word structure that was optimized for numeric processing, coding for it was specialized. The 60-bit word structure of the CDC 6600 and the subsequent Cyber series provided significantly higher precision for floating-point calculations compared to contemporary 32-bit and 36-bit systems. High precision was essential for the recursive calculations involved in gravity modeling. Schutz was an expert in this 60-bit architecture, allowing his team to maximize the hardware’s capability, and that of the locally developed UT-2D operating system.

As geopotential expansions grew more complex and tracking data multiplied, the Computation Center upgraded to systems like the Cyber 170/750 and 175, maintaining compatibility with the old 6600 code but introducing magnetic-core and semiconductor memory enhancements to support operational runs on much larger scales. The Cyber 170/750 was a critical high-performance resource managed by the UT Computation Center. While the university's primary, general-access academic mainframes such as the original CDC 6600 were famously housed in an underground facility beneath the East Mall on the main campus, the Cyber 170/750 was deployed differently. The university began using facilities at the Balcones Research Center, now known as the Pickle Research Campus, to host high-performance systems like the Cyber 170/750. At this off-campus site, the machines were dedicated to handling specialized research and mathematically intensive data analysis tasks. Researchers also utilized a CDC Cyber 175 during the late seventies and eighties. At the time, the Cyber 175 was one of Control Data Corporation’s top-tier, high-speed scalar processors, and it was used intensively for complex finite element and alternating-direction method simulations to solve the types of convection-diffusion problems that frequently appear in reservoir engineering and geology.

In 1985, UT Austin acquired a Cray X-MP, making it the first university in Texas to host a Cray supercomputer. This marked a fundamental paradigm shift from the scalar processing of older mainframes to shared-memory parallel vector computing. The 105 MHz clock cycle and vector pipelining of the Cray X-MP allowed researchers to execute parallelized numerical integrations and invert massive, dense matrices simultaneously. By leveraging the vector capabilities of the Cray X-MP, Schutz pushed the boundaries of gravity modeling even further by running high-precision numerical simulations for the proposed Geopotential Research Mission that modeled the gravity field up to degree and order 180, and even up to degree 300 for certain low orders. Because the Cray used pipelined functional units to execute instructions on entire arrays of data simultaneously, traditional algorithms designed for the CDC and Cyber systems had to be completely restructured to exploit these new vector-parallel capabilities.

Schutz's work directly drove the evolution of the UT Computation Center hardware. The computational system he helped develop to ingest heterogeneous ground and satellite tracking data was known as Utopia, the University of Texas Orbit Processor, first deployed in 1974. Utopia acted as a productionized ingestion engine that merged the numerical integration of orbital equations of motion with a batch least-squares differential correction procedure to seamlessly estimate both satellite orbits and geodetic parameters. Early on, Utopia was used to process satellite laser ranging baseline measurements for the Beacon Explorer mission, and it rapidly adapted to ingest doppler and radar altimetry data from the Geos mission.

\section{McDonald Laser Ranging, Harlan Smith, And Pete Shelus}

During the seventies, the types of measurements flowing into gravity models evolved towards a sophisticated multi-instrumental approach that began to provide the first truly global high-resolution maps of the Earth’s gravitational field. The introduction of satellite radar altimetry and the advent of satellite laser ranging provided powerful new data streams. Collectively, these technologies allowed researchers to fully move beyond simple spherical approximations and begin modeling the geoid, the equipotential surface of the Earth’s gravity field, with unprecedented precision. The foundational method remained the analysis of orbital perturbations, a technique that treats the satellite itself as a test mass in a global physics experiment. As a satellite traverses its orbit, it is subjected to minute accelerations and decelerations caused by the uneven distribution of mass within the Earth’s crust and mantle. During the seventies, these fluctuations were primarily monitored using optical camera networks and the doppler radio-based Transit network. While satellite laser ranging and radar altimetry were the emerging technologies, optical cameras and doppler radio networks provided the sheer volume and global coverage necessary to build the foundational gravity models of that era.

Optical tracking was primarily conducted using Baker-Nunn cameras operated by the Smithsonian Astrophysical Observatory. By capturing a satellite's position against a background of fixed stars, these cameras provided directional measurements with an accuracy of about 10 to 30 meters, or roughly 2 arcseconds. While this was far less precise than the 5-centimeter to 1-meter accuracy of early laser ranging systems, optical networks had a massive structural advantage in global distribution and satellite diversity. For example, during the development of the 1977 Goddard Earth Model 9, 27 different optical stations supplied tracking data for 24 satellites, whereas the new laser systems only tracked four satellites from ten stations. Because these optical cameras captured data across a wide breadth of orbital inclinations from a truly global network, they successfully decoupled highly correlated gravity coefficients and were actually considered the most valuable contributors to early models like GEM-9.

Concurrently, the doppler radio-based networks, such as the Tranet system, tracked the precise shift in radio frequencies to measure a satellite's range-rate, the speed at which the satellite was moving toward or away from the station, with an accuracy of 3 to 7 meters. The Applied Physics Laboratory at Johns Hopkins University pioneered the use of this doppler tracking data, specifically using the Transit navigation satellites, to construct some of the earliest geopotential models, such as APL 1.0 complete to degree 4, and APL 3.5 complete to degree 8.

The broad global structure provided by these two legacy systems was mathematically merged with the emerging high-precision technologies. The milestone Goddard Earth Models such as GEM-9 and GEM-10 synthesized optical directional data, doppler range-rate data, and precise laser ranging data into unified solutions. It wasn't until the launch of dedicated, high-altitude geodetic satellites like Lageos in 1976, combined with the deployment of advanced pulse-chopping laser systems, that laser ranging finally surpassed the optical networks. While the optical and doppler systems successfully defined the broad, global structure of the gravity field through their extensive spatial coverage, the laser network eventually provided the high-precision anchors that defined the future of physical geodesy.

Satellite radar altimetry was a completely new type of technology made possible by the technological advances of the sixties. It was based on directing a radar pulse vertically downward to the ocean surface and measuring the return time. Because the surface of the ocean, when corrected for tides and waves, conforms almost exactly to the geoid, these satellites effectively mapped the shape of the gravity field over the seventy percent of the Earth covered by water. These missions revealed for the first time the gravitational signatures of submerged features like seamounts, trenches, and mid-ocean ridges. The convergence of altimetry data with doppler and laser ranging data allowed for the creation of unified gravity models that bridged the gap between land and sea, laying the technological groundwork for high-fidelity geophysical monitoring systems.

To refine these models and resolve finer details, the seventies saw the full-scale deployment of laser ranging. It created a step-change upward in orbit determination quality by providing precise, unambiguous measurements of the travel time of short laser pulses from ground stations to satellites equipped with specially designed retroreflectors. Since its inception in 1964, laser ranging has evolved from a coarse tracking tool with an accuracy of a few meters into a highly precise system that has transformed understanding of Earth's dynamics and fundamental physics. Schutz and Tapley were involved from the beginning and their research provided the first laser ranging based determination of the Earth's polar motion. By tracking the precise locations of ground stations, laser ranging provided the first space-based empirical validation of geological theories, proving definitively that the Earth's tectonic plates move at rates of a few centimeters per year. It also provided definitive measurements for the Earth's mass and tracked changes in the length of a day and polar motion. Eventually the technology advanced to millimeter-level precision and the effects of General Relativity became observable and had to be factored into orbit models. 

A major reason that laser ranging allowed researchers to model the geoid with such precision was its physical, optical advantage over radio-wavelength tracking systems. At radio frequencies, the permanent electric dipole moment of water vapor in the atmosphere causes molecules to rotate and align with oscillating electric fields, creating a highly volatile wet delay that can cause tens of centimeters of unpredictable error. At the high frequencies of optical lasers, the electric field oscillates too rapidly for heavy water vapor molecules to physically rotate, completely eliminating this dominant delay factor. Because the wet delay in laser ranging is typically only a few millimeters, atmospheric refraction could be corrected using highly stable physical models based purely on ground-station atmospheric pressure. This allowed advanced computational engines, such as the Utopia software developed at the UT Austin, to numerically integrate satellite equations of motion without having to estimate unpredictable dynamic atmospheric parameters, resulting in sub-centimeter orbital tracking consistency. By 1980, third-generation laser systems could operate with compressed pulse lengths of 0.2 nanoseconds, producing single-shot accuracies of seven to fifteen centimeters.

The data from Lageos, in particular, was instrumental in stabilizing the terrestrial reference frame and accurately determining the Earth’s flattening and other low-degree spherical harmonic coefficients, which represent the planet's largest mass irregularities. Launched in May 1976, Lageos was essentially a passive cannonball. It was a dense, sixty centimeter sphere weighing approximately 410 kilograms and covered with 426 retroreflectors. Its unique physical design and orbital placement were intentionally engineered to stabilize the terrestrial reference frame and isolate the Earth's largest mass irregularities. It embodied the ideal characteristics for a certain type of gravity mapping satellite and is an example of uncompromising dedication to performing a very specific role, in this case to act as an ideal test particle free-falling endlessly in the Earth’s gravity field, responding to variations in the gravitational force while minimizing the effects of other forces such as atmospheric drag and solar radiation. It was an ideal target for laser ranging, creating a core component of the gravity modeling workflow taking shape at UT, with measurements flowing from the laser ranging station at UT’s McDonald Observatory back to Austin where the computational resources of the UT Computation Center processed the incoming data in a full-scale operational and productionized environment. 

Because gravitational perturbations naturally attenuate with altitude, Lageos was placed in a very high orbit, with a semi-major axis of about 12,265 km and a mean altitude of roughly 5,900 km, six times higher than most satellites, where it should continue peacefully circling the Earth for at least eight million years. At this height, the satellite was shielded from the noise of short-wavelength gravity variations. This allowed it to act as a pristine gravitational test mass that strongly and uniquely sensed the lowest-degree spherical harmonic terms of the geopotential field, such as the Earth's flattening. By securely defining these broad, low-degree terms, Lageos tracking data allowed geodesists to decorrelate the higher-degree harmonics when its data was combined with observations of lower-orbiting satellites like Starlette, which orbited at roughly 1,100 km and was much more sensitive to finer gravitational structures.

To act as a perfect geodetic anchor, a satellite's motion must be dictated almost entirely by gravity. Because of its extremely low area-to-mass ratio, Lageos was insulated from non-gravitational perturbations that push other satellites off course, such as atmospheric drag and solar radiation pressure. This made its orbital trajectory highly predictable over long periods. Lageos was explicitly tasked with tough geodetic objectives: determining ground station locations to within 10 centimeters and measuring relative tectonic plate motions down to 1 cm/year. Early short-arc analyses using Lageos allowed European stations to compute the 600-kilometer Kootwijk-Wettzell baseline with an accuracy better than 0.7 meters. This ability to precisely locate ground stations established a consistent, geocentric coordinate system. The stable terrestrial reference frame became the essential foundation for orienting global geoid maps and properly locating other geodetic satellites, such as the ocean-mapping radar altimeters. As laser ranging systems and models continued to mature, tracking of Lageos eventually became so precise that it allowed researchers to monitor geocenter motion, detecting the millimeter-level translational movement of the Earth's crustal frame relative to its center of mass, which is caused by the seasonal mass redistribution of global oceans and the atmosphere.

The story of laser ranging at McDonald Observatory and its role in Austin at CSR can be traced back to Harlan Smith. He began his career teaching at Yale University from 1953 to 1963. In 1963, Smith moved to UT Austin to simultaneously become the chair of the Astronomy Department and the director of the McDonald Observatory. He served as department chair for fifteen years and directed the observatory for twenty-six years until his retirement in 1989. As a researcher, Smith's work spanned planetary radio emission analysis, photometry, and instrumentation. 

Beyond his research duties, Smith was a passionate educator and advocate for public outreach. He taught thousands of undergraduate students, developed the award-winning educational film series The Story of the Universe, and created the internationally broadcast Stardate radio program to bring astronomy to the general public. A member of the Quaker faith, he deeply believed in world peace and international scientific cooperation, frequently building bridges by hosting scholars from countries such as China, India, and the Soviet Union. Smith was a passionate advocate for long-term space exploration, strongly supporting the future establishment of lunar bases equipped with astronomical observatories. In his outreach effort with groups like the L5 Society, he played a role during the seventies similar to other public-minded scientific figures such as Carl Sagan and Gerard O’Neill.

Smith was also a pioneering administrator who foresaw the critical role that ground-based observatories would play in supporting space exploration. He successfully convinced NASA, the National Science Foundation, and UT Austin to jointly fund the 107-inch telescope at McDonald Observatory in the sixties. This telescope became the primary ground-based facility for the pioneering lunar laser ranging experiments and was classed as a component in the overall Apollo program. He also served on major national committees, including the National Academy of Sciences committee that helped lay the groundwork for the Hubble Space Telescope, and the NASA Space Science Board, where he played a key role in proposing the Great Observatories series of orbiting telescopes. 

The Hubble became an excellent example of how McDonald integrated into the space age. One specific example of this was the structural role McDonald played in Hubble astrometry. This is especially interesting because of the similarities to McDonald’s role in geodesy. Both astrometry and geodesy are concerned with fundamental reference measurements, of the celestial and terrestrial realms respectively. For Hubble, McDonald provided the highly precise ground-based tracking data necessary to calibrate the spacecraft's instruments. This was achieved through a direct integration between the Texas Minor Planet Program and the Hubble astrometry team.

The primary astrometric instruments aboard Hubble were the fine guidance sensors. To achieve sub-millisecond of arc precision, the astrometry team had to map and correct severe optical field angle distortions caused by the telescope's mirrors and refractive optics. While static distortions were mapped by observing distant, fixed star clusters like M35, the team also needed to perform dynamic calibrations, such as tracking time-varying mechanical flexures, scale changes, and the alignment between the three independent fine guidance units. To do this, they needed to track moving celestial targets with highly predictable paths. Minor planets, or asteroids, were selected as the ideal moving targets.

For the fine guidance calibration to work, the orbital paths of these minor planets had to be known to a higher degree of accuracy than the guidance sensor measurement precision itself. To achieve this, Paul Hemenway and Raynor Duncombe initiated the Texas Minor Planet Program in 1978. Using the long-focus Cassegrain focal plane of the 82-inch telescope at McDonald, the team targeted thirty-four specific minor planets, twenty of which were suited for HST observation. Over the course of the project, researchers captured more than a thousand glass photographic plates of these asteroids. An innovation of the program was the crossing-point method. By taking observations exactly where the projected paths of two different minor planets intersected, researchers could compare the objects against the same local background stars at different times, effectively isolating and eliminating systematic star catalog errors from their orbit reductions.
The precise ground-based observations from McDonald were processed by orbit determination specialists at UT Austin and CSR. By feeding the McDonald data into complex dynamical models, the team refined the asteroid orbits so that their orbital uncertainty was reduced to near-zero. Because these ground-truth reference trajectories were so precise, any deviations or residuals recorded by the fine guidance sensors could be attributed directly to their own instrumental distortions rather than the asteroid's actual motion.

This integration of McDonald Observatory's minor planet data and Hubble astrometry was used to resolve broader issues in celestial mechanics. The combined data allowed researchers to measure the unknown rotation of the European Hipparcos satellite's instrumental coordinate system, effectively linking the stellar optical reference system to a dynamical, gravity-based reference frame. These results were directly incorporated into the celestial reference frame, much as the McDonald geodetic measurements were incorporated in the terrestrial reference frame. Both of these research fields, astrometry and geodesy, can be traced directly to the inspiration from Harlan Smith to focus on participating as fully as possible in NASA’s space missions, tying McDonald and the West Texas Chihuahuan Desert together with the realm of satellites and spacecraft.

To maximize the scientific potential of the new 107-inch telescope and handle the technical challenges of the lunar laser ranging program, Smith recruited specialized experts to the observatory. As part of this drive, he hired Peter Shelus in 1970 specifically to manage, filter, and analyze the complex and noisy photon return data generated by the laser ranging experiments. Equipped with a special Korad ruby laser, the 107-inch became the primary ground-based facility capable of recording regular laser returns from the retroreflector arrays deployed by the Apollo astronauts. For years, it operated as the primary lunar ranging program in the world. 

In the early days of the Apollo program, conventional astronomy and lunar laser ranging shared time on the 107-inch, which limited lunar observations to three 45-minute sessions per day and carried high operational costs. Shelus led the transition of the program into a dedicated, highly automated, dual-purpose station known as the McDonald laser ranging station. Utilizing a 0.76-meter Cassegrain reflecting telescope eventually situated on Mount Fowlkes, the MLRS was designed for routine ranging of both artificial satellites and the Moon. By migrating the ranging operations to the automated MLRS, the observatory lowered long-term operational costs, increased the data acquisition rate for ranging, and freed the 107-inch telescope to return to conventional uses. 

Because the return signal strength from the Moon scales with the inverse fourth power of the distance, a laser pulse might yield fewer than ten returning photons. To isolate these rare signals from solar background noise and dark counts from the detector, Shelus directed the implementation of a rigorous four-stage filtering system at the MLRS. The photodetector's gate was opened during a narrow time window centered on the expected 2.5-second round-trip travel time. Utilizing a pinhole aperture in the focal plane limited the detector's field of view to a tiny portion of the lunar surface. High-rejection optical bandpass filters blocked all incoming light except the specific 532 nm green wavelength of the station's laser. And timing histograms were analyzed with statistical algorithms to separate random background noise from the coherent peak of the retroreflected pulse. Shelus’s research also advanced mathematical data-filtering techniques to successfully extract small signal returns from high background noise.

Shelus served as a crucial bridge between the laser ranging and global geodetic communities, He worked to ensure that the high-precision tracking data generated at the McDonald Observatory was seamlessly integrated into the world's foundational reference frames, and made significant contributions to the International Earth Rotation and Reference Systems Service, the International Terrestrial Reference Frame, and the dynamical realization of the celestial reference frame. He was deeply embedded in the administrative and scientific infrastructure of the IERS. He served as the director of the lunar laser ranging coordinating center for the IERS and represented the community on the IERS directing board.

Long before the modern IERS was fully established, Shelus managed McDonald Observatory's participation in the Earth Rotation by Lunar Distances campaign during the late seventies. He directed the filtering, compression, and transmission of McDonald lunar laser ranging data to the Bureau International de l'Heure in Paris, enabling the extraction of Earth rotation and polar motion parameters. As a member of the IERS directing board, he played a key role during the organization's transitional periods. For example, he represented the laser ranging community in discussions regarding the implementation of International Astronomical Union resolutions within the IERS, governing Earth precession and nutation models, the definition of Terrestrial Time, and transformations between terrestrial and celestial systems. This work effectively carried on Ray Duncombe’s legacy, stretching back before the dawn of the space age.

These efforts were instrumental in the realization of the International Terrestrial Reference Frame (ITRF), linking satellite motion to the physical dynamics of the Earth’s crust. The ITRF and International Celestial Reference Frame (ICRF) are fundamental coordinate systems used to map the Earth and outer-space, respectively. To accurately model satellite orbits and Earth's dynamics, it is essential to track the complex mathematical transformations between the Earth-fixed ITRF and the space-fixed ICRF. Because the Earth’s shape is constantly deforming due to tides and tectonic plate motion, a truly fixed Earth frame doesn't naturally exist. Instead, the ITRF is realized through a global ground network of tracking stations utilizing laser ranging, doppler, and more recently GPS. 

The terrestrial frame provides the highly accurate, standard coordinates and velocities for a network of dedicated ground stations, accounting for continental drift and crustal deformation. It relies on the continuous monitoring of Earth orientation parameters such as the Earth's variable rotation rate, polar motion, and variation of latitude. Shelus co-authored numerous analyses that processed decades of McDonald laser ranging data to estimate these parameters. Under his direction, the McDonald laser ranging operations center produced orientation values and computed highly precise geocentric station coordinates which were provided directly to the IERS Terrestrial Frame Section to track the orientation of the global reference frame. Shelus focused the operations center on securing high-quality observations near the new and full moon phases. Because lunar laser ranging is heavily dominated by the lunar month cycle, fixing these data gaps was critical for accurately decoupling true geophysical and Earth-rotation signals from systematic and relativistic effects.

The celestial frame is a space-fixed, quasi-inertial reference frame. Rather than relying on the dynamical motion of planets as in the past, the modern ICRF is a kinematic frame defined entirely by observations of distant, extragalactic radio sources whose proper motion is negligible. These quasars are observed and measured primarily using Very Long Baseline Interferometry. Meanwhile, lunar laser ranging contributes to the definition of a dynamically realized inertial reference frame based on the equations of motion of the Earth-Moon system. Shelus’s work maintaining the long-term, highly stable laser ranging database provided the necessary data to define this dynamical reference frame to an accuracy of milli-arcseconds. Comparing this gravity-based dynamical frame against the kinematically realized ICRF is essential for mutual cross-checking, identifying systematic biases in the optical and radio reference systems, and validating the overall stability of cosmic distance and coordinate scales. It has been instrumental in testing gravitational physics, such as setting tight constraints on variations in the gravitational constant, and in realizing dynamically defined terrestrial reference frames and orientation.

McDonald’s laser ranging system transitioned into a dual-purpose facility that tracked both the Moon and a growing roster of artificial satellites. While his early career heavily emphasized lunar laser ranging, under his leadership as Operations Center Manager in the mid-eighties, the observing emphasis at MLRS shifted dramatically toward satellite laser ranging. To execute this dual-purpose mission, the MLRS's 0.76-meter telescope was equipped with a transmitter design that integrated two distinct lasers sharing a common control and power electronics backbone. For ranging to artificial satellites, the station employed a specialized laser, allowing the station to successfully target fast-moving, near-Earth objects.

Shelus frequently highlighted the immense physical contrast between ranging artificial satellites and ranging the Moon. Because signal strength drops off based on the inverse fourth power of the target's distance, laser ranging to the Moon is more than a trillion times more difficult than ranging to near-Earth artificial satellites. While satellite signals were much stronger, his team had to adapt to tracking targets with vastly different apparent angular speeds across the sky. Low-orbit satellites move very quickly compared to the slow-moving Moon or high-altitude satellites. The MLRS was designed to be versatile and less demanding for the observing crew. During a single 8-hour shift, an operator could seamlessly pivot between the Moon and a diverse suite of artificial satellites. The routine satellite laser ranging target list at McDonald grew to include major geodetic and oceanographic missions such as Lageos, Starlette, and Topex. In collaboration with the NASA crustal dynamics project, the station integrated workstations to handle the high volume of satellite data. This allowed the team to automate the filtering of satellite laser ranging data on-site.

Because it served as a highly stable, fixed fiducial point on the North American tectonic plate, the observatory provided a thirty-plus year record of geodetic positioning data that remains a cornerstone for the realization of the modern terrestrial frame. Satellite laser ranging is uniquely responsible for defining the origin of the Earth's exact center of mass or geocenter, and monitoring its movement over time. It also determines the Earth's mass and gravitational constant. It is privileged for defining the terrestrial origin because its translation components show less scattered temporal behavior compared to other techniques like doppler or gps. Because the Earth wobbles and its rotation speed varies, it is essential to measure orientation parameters like polar motion and length of day. These align the terrestrial and celestial frames, and provide a high-precision time-series of orientation parameters, making it possible to interpolate the Earth's wobble and rotation measurements between radio astronomy campaigns.

Research progress on all of these diverse fronts flowed back from McDonald to Austin, where they were integrated and processed along with a variety of other, independent data sources. This was the central role of CSR and the numerical code developed by Bob Schutz on the CDC 6600 and its Cyber descendants. While Ray Duncombe and Harlan Smith had established the foundations and worked tirelessly in the sixties to clear the way, it was researchers such as Schutz and Shelus that brought the fully operational system into being during the seventies.

\section{UTOPIA}

The new era of satellite altimetry and laser ranging was embodied in the development of Utopia, the University of Texas Orbit Processor. With a variety of data streams flowing in, and the push towards more detailed gravity maps, the scale of the problem was large and growing both in terms of the spherical harmonic representations and the measurement data sides of the gravity models. From a purely computing perspective, the push for higher-resolution gravity maps required expanding spherical harmonic models to higher degrees and orders, which exponentially increased the sizes of the matrix equations. In the early days, gravity models were limited to degree and order eight or sixteen. By the time UT Austin was processing data for the Topex mission, models like the joint gravity model were expanded to degree and order seventy. This required Utopia to simultaneously estimate over 5,000 to 7,000 distinct gravity and orbital parameters, turning the orbit processor into a massive engine for solving thousands of coupled linear equations.

Likewise for the measurement data, particularly for the square covariance matrices representing the uncertainties in the measurements. Since the measurements should be weighted by their uncertainties, much like a weighted average, these uncertainty matrices play a critical role. Factorizations to handle massive square matrices played an important role as simpler least squares methods became computationally impractical due to computer round-off and truncation errors. To address this, Schutz and his coworkers implemented advanced divide-and-conquer linear algebra techniques in Utopia, specifically the square-root information filter and smoother and upper-diagonal covariance factorization. Instead of operating on the full covariance matrix directly, these algorithms used orthogonal transformations to break down and factorize the matrices into upper-triangular forms. This effectively doubled the numerical precision of the calculations for a given computer word length, allowing the software to reliably process massive batches of tracking data without the filter diverging or crashing.

There were actually two generations of software, as part of the long-term software efforts led by Tapley and Schutz, beginning with Earthod. It was written in Fortran IV on the 6600 and naturally led onward toward Utopia on the Cyber and Cray. Earthod was primarily used as a simulation program. It generated simulated range, range-rate, elevation, and azimuth measurements for a satellite observed by up to twelve ground-based tracking stations. It was specifically utilized for testing adaptive and online filtering procedures to estimate a satellite's state from measurements corrupted by noise. Since it simulated both the true values and corrupted values, it enabled research into the nature and effects of measurement uncertainties. 

In the mid-seventies, CSR researchers utilized Earthod to study the Transit system by processing simulated doppler range-rate observations to test adaptive filtering procedures in a rectangular coordinate system. These simulations were crucial because they revealed the computational limits of the era's methodologies. While Earthod's localized adaptive filters successfully reduced positional uncertainty by about four meters, they only improved absolute positional accuracy by 0.6 meters. Earthod effectively demonstrated that localized, short-arc rectangular filters were structurally inadequate for resolving global systematic errors or separating radial, along-track, and cross-track errors over single satellite passes. It served as the necessary controlled environment to master the effects of noise, filtering, and coordinate frame mismatches. By successfully mapping how these uncertainties propagated in simulation, Schutz and his coworkers ensured that when the transition to Utopia occurred, the software possessed the robust statistical architecture required to validate trajectories for landmark missions like Lageos, Seasat, and Topex.

The limitations discovered in Earthod necessitated the development of Utopia, a robust, operational orbit processor. While Earthod was capped at twelve stations and relied on localized filters and analytical approximations, Utopia was designed to handle unlimited global tracking networks, moving beyond simple range and azimuth simulations to process massive batches of real-world satellite laser ranging, altimetry, doppler, and gps data. Instead of short-arc rectangular filtering, Utopia employed a global estimator with a highly parameterized dynamic model, utilizing both batch least squares and extended sequential filters such as the square-root information filter. The computational power of the Cyber and Cray platforms allowed Utopia to integrate complex differential equations over multi-week arcs, estimating thousands of parameters simultaneously, including high-degree spherical harmonics, ocean tides, and empirical non-conservative forces.

Unlike Earthod, which had been used primarily in the context of simulation and filtering experiments using the CDC 6600, Utopia was used operationally to compute and integrate highly precise satellite trajectories. It was also used to validate orbital data processed by NASA's own orbital analysis software during highly precise tests of fundamental physics. The specialized Lageos satellites were used to validate general relativity's frame-dragging effect. By using Utopia to rigorously track the nodes of these counter-orbiting butterfly configuration satellites over an eleven year period, researchers successfully isolated the minute force of the rotating Earth dragging spacetime around with it. Because these tests of fundamental physics required measuring orbital perturbations at the level of milliarcseconds, the scientific community could not rely on a single mathematical model. Utopia served as a vital, independent verification engine. The data was processed through NASA's Geodyn software, the German Epos software, and UT Austin's Utopia. By comparing the results from these completely distinct digital infrastructures, researchers could prove that the relativistic signals they were seeing were actual physical phenomena, and not just systematic biases or artifacts hidden in a particular organization’s code.

As tracking networks expanded, researchers quickly realized that the problem of orbit determination had become heavily over-determined, meaning there were far more observations than parameters to solve for, and all of them were corrupted by noise. This transition quickly exposed the limitations of purely deterministic approaches when confronted with the imperfect observations of global tracking networks, and imperfect knowledge of the forces involved in the differential equations governing celestial motion. Utopia introduced statistical estimation utilizing least squares methods, expressed as large-scale linear algebra problems coded in Fortran IV. It was specifically engineered to tackle these matrix equations by implementing two statistical estimation engines. First, a classical weighted batch least-squares estimator to process large historical datasets offline. And second, an extended sequential filter to process data sequentially. This allowed the software to estimate a satellite's precise state by minimizing the squared errors of imperfect tracking data while simultaneously absorbing the unmodeled forces of the environment. The objective was routine computation of empirically accurate orbits and gravity models from masses of real-world measurement data flowing in from ground-stations and altimetry satellites. 

For context on the historical role of Utopia, it is useful to look at the related Geodyn project developed by NASA’s Goddard Space Flight Center. Utopia and Geodyn were both leading orbit determination software systems that matured in the seventies and were utilized by the global community. They were frequently used in parallel by researchers to independently verify complex measurements, but they had distinct architectural and mathematical differences. While Utopia was born on the CDC 6600 and optimized for the Cyber and Cray supercomputers, Geodyn was historically tied to NASA's IBM mainframes such as the IBM 7094 and 360 before migrating to platforms like the Cyber 205. The CDC Cyber 205 was a famous vector-processor supercomputer released in 1980. During the early-to-mid eighties, the Cyber 205 battled fiercely with the Cray-1, and later the Cray X-MP, for the title of the fastest computer in the world.

To propagate orbits, Geodyn specifically relied on a predictor-corrector method, contrasting with Utopia's fixed-step, fixed-order integrator. Geodyn utilized a traditional batch least squares approach for its parameter estimation, while Utopia implemented a square-root information filter and smoother, although it is also capable of weighted least squares batch procedures. Geodyn had built-in capabilities to handle multi-arc, multi-satellite processing, while Utopia was designed specifically as a single-satellite program. To handle multi-satellite scenarios, CSR developed a separate but related multi-satellite orbit determination program. 

The introduction of laser ranging and radar altimetry greatly increased the scale of the incoming measurement dataflows and the need for floating-point processing power. The scale of incoming data transformed the discipline very quickly. Earlier tracking relied heavily on optical cameras or early radio interferometers, which provided relatively sparse observations. The advent of continuous measurement systems, especially the radar altimeters on Geos in 1975 and Seasat in 1978, alongside the rapid expansion of the global satellite laser ranging network, created an unprecedented flow of data. To process Seasat data and map the marine geoid, Utopia had to ingest and time-sample altimeter measurements at twenty second intervals across multi-week orbital arcs.

The physical limitations of the CDC hardware directly informed Utopia’s numerical code, as the need to process global observations surpassed the raw storage capacity of the machines, requiring a nuanced approach to numerical linear algebra to handle the increasing volume of satellite laser ranging and altimetry data. There was a strict memory bottleneck on the CDC 6000 and Cyber 170 series. Because the architecture relied on eighteen bit addressing, the systems were subject to a hard limit of only 262,144 sixty bit words of central memory. While the machines could be supplemented with extended core storage that held millions of words, it could only be used for high-speed block swapping and data storage.
Rather than relying on a brute force approach, naively using random-access memory, which would have been completely impractical on seventies hardware, Utopia emphasized sparse matrix representations and a preference for sequential, filtering updates for large-scale gravity models, along with the state transition matrix concept to relate differential changes in initial conditions to observed residuals across time intervals. And it incorporated square root computational algorithms, with matrix factorizations and transformations to operate directly on information arrays, rather than the more numerically challenging covariances. 

Because orbital dynamics are governed by highly complex, nonlinear differential equations, attempting to compute the satellite's exact state for every single tracking observation across a multi-week arc would have crippled the system. Instead, Utopia utilized the state transition matrix as a linear mapping tool, allowing the system to linearly map deviations in the state vector from an initial reference epoch to any subsequent observation time. This elegantly related any differential changes in the initial conditions directly to the observed tracking residuals, bypassing the need to continuously re-solve the nonlinear trajectory from scratch for each data point.

Because the true trajectory of a satellite is never known exactly, Utopia first numerically integrated a nominal reference trajectory based on a best-guess initial state. By expanding the highly non-linear differential equations of motion into a series about this reference trajectory, the software replaced the original complex non-linear problem with a much simpler linear estimation problem focused purely on calculating state deviations, the difference between the true state and the reference state. The state transition matrix represented the linear term of this series expansion. To generate this matrix, Utopia numerically integrated the variational equations of motion along the reference trajectory, using the matrix of partial derivatives of the force model evaluated on the reference orbit. 

The core computational challenge was to relate an observation made at some later time back to the initial conditions. Utopia achieved this via a measurement sensitivity matrix, leading to a mapped observation matrix. The result linearly linked the observed tracking residuals at any arbitrary time back to the unknown initial state deviation, while accounting for measurement noise. Utopia could then sequentially read thousands of laser ranging or altimetry observations across a multi-week arc and mathematically map each of them back to a single reference epoch. By expressing all observations in terms of the state at a single epoch, it reduced the number of unknown state vectors from one-per-observation down to a single set of initial conditions, preventing the system from having to constantly re-integrate the non-linear trajectory from scratch for each data point.

Utopia favored operating on information arrays rather than traditional covariance matrices. Standard sequential estimators directly updated the covariance matrix, but with the finite-digit arithmetic of seventies computers, these matrices would often succumb to truncation and rounding errors. This was especially a problem as highly accurate laser ranging data was introduced. Because of numerical disparities and mismatches, the mathematical symmetry and positive-definite character of covariance matrices could quickly degrade, causing the filter to diverge or crash. By operating on the information matrix instead, the filter accumulated information sequentially without suffering from these catastrophic numerical instabilities. And with a square-root information filter, the code operates on the square root of the information matrix, reducing the condition number of the matrices being manipulated. A matrix's sensitivity to computer rounding errors is determined by its condition number. By mathematically operating on the square root of the matrix, the condition number was effectively halved. This essentially doubled the numerical precision of the computer for a given word length and ensured that the matrix remained healthy.

Because of its inherent stability, the square-root information filter was exceptionally well-suited for processing the large volumes of highly accurate observations becoming available in the seventies. It allowed Utopia to sequentially process very large amounts of tracking data over long arcs without the accumulated computer precision errors that would crash standard batch processors or traditional filters. Utopia could be initialized with infinite or singular covariance matrices, meaning zero initial confidence, and still yield highly accurate results once enough data was processed. This ensured that the centimeter-level accuracy inherent in the raw tracking data was not lost to computational errors during the filtering process. 

While the square-root approach solved the stability problem, there was a secondary hardware issue. Calculating actual mathematical square roots was highly taxing for seventies hardware. To save computational overhead, developers introduced square-root free equivalents. This led to the use of the upper-diagonal covariance factorization, which decomposed the matrix into a unit upper triangular matrix and a diagonal matrix. By manipulating these components separately using orthogonal transformations, in a sense converting to a geometric problem, Utopia achieved all the numerical stability of a square-root filter without forcing the hardware to actually calculate computationally expensive square roots. An orthogonal transformation preserves the length or norm of vectors while rearranging the matrix into a solvable upper-triangular form. One approach zeroes out elements below the main diagonal one at a time. By calculating specific sine and cosine rotation parameters, it iteratively eliminates targeted data points. This element-by-element approach makes it highly flexible for sequentially processing incoming streams of satellite observations. Instead of rotating a plane, another method uses an elementary reflector, essentially treating a mathematical plane as a mirror. The primary advantage is its efficiency. It can zero out an entire column of elements below the main diagonal in a single computational operation, rather than eliminating them one by one.

Another cause of numerical stability issues in orbit determination has to do with positional coordinates near the Earth’s poles, especially for satellites in polar orbits, a common situation in gravity modeling. The root cause here is the coordinate representation used, with common spherical representations such as latitude and longitude effectively breaking down at the poles. At the north pole, longitude becomes indeterminate and latitude reaches its maximum value of ninety degrees. For gravity field models using spherical harmonics, as the models grew in complexity to include higher degrees and orders, the traditional equations suffered from mathematical singularities and numerical instabilities as longitude effectively broke-down near the poles. 

In 1973, a new formulation of spherical harmonics successfully removed the problematic latitude and longitude parameters, replacing them with dimensionless direction cosines in a rectangular Earth-fixed rotating frame to bypass the division-by-zero error at the poles. These equations relied on unnormalized functions however, and as researchers attempted to model the gravity field at higher resolutions, the terms in this unnormalized series grew exponentially. For anything beyond the fourth degree and order, or for integration arcs longer than about three hours, these massive numbers quickly exceeded the finite floating-point limits of seventies computers, causing numerical overflows and underflows. While the original rectangular formulation was adequate for short integrations, it needed to be expanded for more rigorous analysis. 

To solve this scaling crisis, Schutz and collaborators developed a fully normalized variant of the spherical harmonics representation. By introducing a normalization factor to safely scale the numerical calculations, they eliminated overflows and underflows, keeping the results bounded and within the computer's precision limits regardless of how high the degree and order climbed. This normalization is what unlocked Utopia's scalability. For Lageos, implementing the normalized formulation allowed Utopia to safely extend its orbit integrations from three hours to twenty-four hours, successfully processing a gravity model complete to the seventh degree and order, and capturing minute perturbations acting on the satellite. For Seasat and Topex, Utopia was able to handle massive gravity fields, integrating datasets like satellite-to-satellite tracking and radar altimetry crossovers. 

Utopia eventually scaled to much higher resolutions, up to 70 by 70 and 180 by 180, as computational power increased. The methodology was explicitly designed to maintain both computational speed and numerical stability across all latitudes. This made it the ideal mathematical engine for processing the large batches of tracking data required for precision gravity modeling and modern geodynamics. The rectangular coordinates approach championed by Schutz proved so robust that it became a foundational strategy in astrodynamics, and was later adopted to solve similar polar breakdown issues in other geodetic domains, such as evaluating global atmospheric density for low-Earth orbiters, which originally suffered from the same longitudinal singularities at the poles.

Utopia’s overall framework integrated complex force models, including Earth tides and solar radiation pressure, while providing a critical tool for assessing the accuracy of the numerical integrations it was performing. This allowed it to process the first generation of high-precision laser ranging data with unprecedented fidelity. While passive, spherical satellites like Lageos and Starlette naturally minimized non-gravitational forces, complex active satellites like Seasat possessed large, asymmetric solar panels and instruments that made their aerodynamic and radiation parameters highly uncertain and virtually unobservable. Rather than expanding the state vector to actively estimate hundreds of specific physical parameters for these forces, Utopia estimated empirical once per revolution along-track and cross-track accelerations over short sub-arcs. This parameterization gracefully absorbed the growth of errors in the orbital node, inclination, and eccentricity caused by mis-modeled surface forces, preventing trajectory degradation.
 
Within this analytical environment, the consider-covariance analysis feature emerged as an important tool for sensitivity analysis, allowing researchers to assess the impact of unestimated parameters on the state estimate and thereby structure more robust estimation algorithms. This is a technique used to account for the uncertainties of parameters that are present in a model but are not actively being estimated. In orbit determination and gravity modeling, there are many parameter values that contain errors, and three basic options for dealing with them. The simplest is to estimate the state while neglecting any errors in the non-state parameters. Or the state can be expanded to actively estimate parameters that might be in error. The third option is that the parameters are not actively estimated, but their known statistical uncertainties are mathematically considered or represented in the uncertainty matrix. Some tradeoffs are that keeping the parameter array as small as possible saves computer memory, processing time, and computational costs. And there are simply limits to what can be done, since large arrays of parameters often cannot be completely observed or uniquely determined from a limited set of tracking data collected over a short time interval, and attempting to estimate unobservable parameters tends to lead to poor results and outright errors.

By running a consider-covariance analysis, it is possible to quantify exactly how much the errors in unestimated parameters such as atmospheric drag, solar radiation, or other incorrect assumptions will degrade the accuracy of the overall state estimate. It is then possible to identify the optimal set of parameters to be actively estimated, versus those that can safely be ignored. In sequential estimation modes such as Utopia’s square root information filter, a danger is that the calculated uncertainties can become too small over time, causing the filter to become overly optimistic and unresponsive. Using a consider-state approach can make the filter more sensitive by enforcing a more realistic uncertainty matrix, resulting in more robust performance. 

Schutz and CSR researchers used the consider-covariance approach together with tracking data from Starlette, a small, dense, spherical satellite that orbited at a lower altitude than Lageos, making its trajectory exceptionally sensitive to the Earth’s gravity field. Utopia was able to evaluate how inaccuracies in the static gravity field impacted Earth rotation parameter calculations for polar motion and length of day, and revealed that errors in first-order geopotential coefficients were inducing a systematic perturbation with a dominant seventy-four day period in the polar motion solutions. They were then able to improve the force model by simultaneously estimating selected geopotential and ocean tide parameters. This robust adjustment significantly reduced the discrepancy between the Starlette and Lageos Earth rotation parameter estimates, demonstrating the power of the approach for isolating and neutralizing complex geophysical errors.

With Utopia operational from the late seventies onward, CSR became a primary analysis center for satellite laser ranging data and began releasing annual tracking station position and velocity solutions and Earth orientation parameters to international bodies. This involved analyzing global network observations of satellites such as Lageos on a weekly basis, providing some of the most stable and longest-spanning polar motion series used in the international Earth rotation service. Utopia's centimeter-level precision relied on integrating the most advanced force models of the era. The software utilized the Joint Gravity Model 3, which combined data from dozens of satellites to resolve the static gravity field complete to spherical harmonic degree and order seventy, drastically reducing geographically correlated orbit errors. To handle the dynamic gravitational pull of water mass, researchers developed the CSR 3.0 global ocean tide model using altimetry data from Topex. By synthesizing these massive gravity and tidal models, Utopia could minimize residual discrepancies and achieve radial orbit accuracies of two to three centimeters for critical altimetry missions.

An example of the Utopia system in operation involved laser ranging data and the Ajisai satellite to estimate first-order geodetic control points. Ajisai’s low orbit made it highly susceptible to errors in atmospheric density and geopotential models, meaning the forces changing the orbit could not be simply modeled over long intervals. To handle this situation, the entire multi-month observation arc was divided into three day sub-arcs. While the satellite's state was adjusted every three days to absorb dynamic errors, global parameters, such as the actual tracking station positions and radiation pressure coefficients, were estimated once globally for the entire arc. This global multi-arc approach successfully achieved a geocentric station coordinate accuracy of three to four centimeters. In contrast, the prior short-arc method, which only used two successive satellite passes simultaneously, lacked global constraints and resulted in absolute rectangular coordinate discrepancies of more than ten centimeters. This milestone facilitated wide-ranging Earth science studies and established the foundation for tracking regional plate motion and crustal deformation with unprecedented reliability. 

While these terrestrial-based measurements defined the Earth’s orientation and surface geometry, a truly unified reference architecture required a dual-pronged approach that extended into the celestial reference frame through lunar distance measurements. Lunar laser ranging at McDonald was the unique technique capable of linking the Earth-fixed terrestrial system to the inertial, space-fixed celestial system. This high-precision analysis was essential for determining the orientation of the dynamical ecliptic and refining the Earth’s precession constant. Furthermore, it provided the most accurate estimation of the tidal acceleration of the Moon, the energy dissipation of the Earth-Moon system, and the increase in the lunar distance of approximately four centimeters per year.

By the late seventies and early eighties, Bob Schutz’s work in bringing the Utopia system to life served as a role-model and archetype for computational modeling and simulation of the era. It was designed from the ground up to run on the CDC 6600, Cyber, and Cray. It had become fairly remote from general purpose computing, and was dedicated to productionized operational workflows of global data and regular, standardized numerical processing. All of this was far from the concerns of general purpose computer users at UT Austin, who at most had possibly heard about the new machines and software indirectly, and even perhaps glimpsed the additional blinking lights of the hardware in the UT Computation Center machine room. This was a change from earlier days, when the CDC 6600 ruled alone and everyone coexisted on effectively the same all-encomposing machine, with numerical code such as Earthod running alongside student software projects, many of which were not even coded in Fortran, and where a parallel environment and community were taking root. 
